\documentclass[12pt,oneside,openany,pagenumber=footcenter]{book}
%\documentclass[a4paper,10pt]{report}

\usepackage[utf8]{inputenc}
\usepackage{lmodern}
\usepackage[T1]{fontenc}
\usepackage{amsfonts}

\usepackage{slovak}
\usepackage{fontenc}
\usepackage{graphicx}
\usepackage{graphics}
\usepackage{graphicx}
\usepackage{hyperref}
\usepackage{makeidx}

\newtheorem{definicia}{Definícia}[section]
\newtheorem{HLPpoznamka}{Poznámka}[section]
\newtheorem{HLPpriklad}{Príklad}[section]
\newtheorem{HLPcvicenie}[HLPpriklad]{Cvičenie}
\newtheorem{zadanie}{Úloha}[section]
\newenvironment{poznamka}{\begin{HLPpoznamka}\rm}{\end{HLPpoznamka}}
\newenvironment{priklad}{\begin{HLPpriklad}\rm}{\end{HLPpriklad}}
\newenvironment{cvicenie}{\begin{HLPcvicenie}\rm}{\end{HLPcvicenie}}
\newtheorem{veta}{Veta}[section]
\newtheorem{lema}[veta]{Lema}
\newtheorem{dosledok}[veta]{Dôsledok}
\newtheorem{teza}[veta]{Téza}
\newtheorem{dokaz}[veta]{Dôkaz}


\pagestyle{headings}

\bibliographystyle{unsrt}

\def\indexname{Register}
% pekne pokope definujeme potrebne udaje
\def\mftitlea{Biologicky motivované výpočtové modely}
\def\mftitle{\mftitlea}
\def\mfthesistype{Dizertačná práca}
\def\mfauthor{Michal Kováč}
\def\mfadvisor{doc. RNDr. Damas Gruska, PhD.}
\def\mfplacedate{Bratislava, 2011}


\ifx\pdfoutput\undefined\relax\else\pdfinfo{ /Title (\mftitle) /Author (\mfauthor) /Creator (PDFLaTeX) } \fi

\def\eps{\varepsilon}
\def\goodgap{\hspace{\subfigcapskip}}
\makeindex

\begin{document}
\frontmatter
\thispagestyle{empty}
\begin{minipage}{0.20\textwidth}
\includegraphics[width=0.9\textwidth]{img/comenius_half.png}
\end{minipage}
\begin{minipage}{0.79\textwidth}
\begin{center}
\sc Katedra Informatiky \\
Fakulta Matematiky, Fyziky a Informatiky \\
Univerzita Komenského, Bratislava
\end{center}
\end{minipage}

\vfill
\begin{center}
\begin{minipage}{0.8\textwidth}
\hrule
\bigskip\bigskip
\centerline{\LARGE\sc\mftitlea}
\smallskip
\centerline{(\mfthesistype)}
\bigskip
\bigskip
\centerline{\large\sc\mfauthor}
\bigskip\bigskip
\hrule
\end{minipage}
\end{center}
\vfill
{\bf Vedúci:} \mfadvisor
\hfill\mfplacedate
\eject
\eject

\thispagestyle{empty}
{~}\vspace{12cm}

{~}\vspace{12cm}

\begin{minipage}{0.25\textwidth}~\end{minipage}
\begin{minipage}{0.69\textwidth}
I hereby declare that this submission is my own work and that, to the best of my knowledge and belief, it contains no material previously published or written by another person nor material which to a substantial extent has been accepted for the award of any other degree or diploma of the university or other institute of higher learning, except where due acknowledgment has been made in the text.

\bigskip\bigskip

\hfill\hbox to 6cm{\dotfill}
\end{minipage}

\chapter*{Acknowledgements}

I am very grateful to all the people I had the privilege to work with, especially I would like to thank my advisor RNDr. Damas Gruska PhD. for his constant interest, encouragement and help that I deeply appreciate.
The present work was possible due to many beneficial consultations and intensive cooperation.

I would also like to thank to the members of Laboratory of Comparative and Functional Genomics of Eukaryotic Organelles and their head, Prof. Ľubomír Tomáška for familiarization with the work of biologists.

I thank to my family, colleagues, teammates and all the friends who supported me especially last weeks before the deadlines for their patience, support and understanding\footnote{Understanding that they will propably never understand my work;)}

% Osobitná vďaka patrí vedúcemu diplomovej práce doc. RNDr. Damasovi Gruskovi PhD. za cenné rady, námety, podnetné pripomienky a všestrannú pomoc, ktorú si hlboko vážim. Len vďaka mnohým prínosným konzultáciam a intenzívnej spolupráci som bol schopný napísať toto dielo. Nesmiem zabudnúť ani na RNDr. Branislava Rovana, CSc. a spolužiakov za to, že si na dizertačnom seminári našli čas, aby si vypočuli moju prezentáciu dizertačnej práce. Ďalšie poďakovania venujem rodičom a známym, ktorí to so mnou dokázali vydržať posledné týždne pred odovzdaním.

% !TEX root = diz.tex
\chapter*{Abstract}
\begin{description} \itemsep1pt \parskip0pt \parsep0pt
  \item[Author:] \mfauthor
  \item[Title:] \mftitle
  \item[University:] Comenius University in Bratislava
  \item[Faculty:] Faculty of Mathematics, Physics and Informatics
  \item[Department:] Department of Applied Informatics
  \item[Supervisor:] \mfadvisor
\end{description}

This work discusses the research in the membrane systems, an emerging field of natural computing. Many variants of membranes systems have already been studied, most of them uses parallel rewriting and are computationally complete. Various sequential models have been proposed, however, in many cases they are weaker than their parallel variant.

We propose some other sequential models, provide universality proofs for them and suggest further research.

\begin{description}
  \item[Keywords:] Computation models inspired by biology, Membrane systems, P systems
\end{description}

\chapter*{Abstrakt}
\begin{description} \itemsep1pt \parskip0pt \parsep0pt
  \item[Autor:] \mfauthor
  \item[Názov dizertačnej práce:] \mftitle
  \item[Škola:] Univerzita Komenského v Bratislave
  \item[Fakulta:] Fakulta matematiky, fyziky a informatiky
  \item[Katedra:] Katedra aplikovanej informatiky
  \item[Vedúci dizertačnej práce:] \mfadvisor
  \item \mfplacedate
\end{description}

V tejto práci sa zaoberáme výskumom v oblasti membránových systémov. Veľa variantov už bolo preskúmaných, väčšinou pri výpočte používajú paralelizmus a sú Turingovsky úplné. Navrhlo sa aj veľa sekvenčných modelov, ale väčšina z nich má slabšiu výpočtovú silu ako ich paralelný variant.

V druhej časti predkladáme niektoré sekvenčné modely, u ktorých dokazujeme Turingovu úplnosť a navrhujeme modely, ktoré plánujeme preskúmať v rámci dizertačnej práce.

\begin{description}
  \item[Kľúčové slová:] Výpočtové modely inšpirované biológiou, Membránové systémy, P systémy
\end{description}
\tableofcontents{}
\listoffigures{}
\listoftables{}

\mainmatter
% !TEX root = diz.tex
\addcontentsline{toc}{chapter}{Introduction}

There are a lot of areas in the theoretical computer science that are motivated by other science fields. Computation models motivated by biology forms a large group of them. They include neural networks, computational models based on DNA evolutionary algorithms, which have already found their use in computer science and proved that it is worth to be inspired by biology. L-systems are specialized for describing the growth of plants, but they have also found the applications in computer graphics, especially in fractal geometry. Other emerging areas are still awaiting for their more significant uses.

One of them is the membrane computing \cite{Paun10OxfordHandbookMembraneComputing}. It is relatively young field of natural computing - in comparison: neural networks have been researched since 1943 and membrane systems since 1998 \cite{Paun98}.

Membrane systems (P systems) are distributed parallel computing devices inspired by the structure and functionality of cells. Recently, many P system variants have been developed in order to simulate the cells more realistically or just to improve the computational power.

We will start by an introduction of some natural computing areas including models inspired by biology in Chapter \ref{cha:natural_computing}. In Chapter \ref{cha:preliminaries} we recall some computer science basic notions that we will use through the work. P systems are formally presented in Chapter \ref{cha:p_systems}, with the current state of the research in their variants, overview of software simulator MeCoSym and various case studies.

In Chapter \ref{cha:on_the_edge_of_universality_of_sequential_p_systems} we will present the current state of our work, mainly from theoretic viewpoint (computational power, decidability of behavioral properties), including the published results in sections \ref{sec:inhibitors}, \ref{sec:active_membranes} and \ref{sec:notions_from_reaction_systems}.


\chapter{Preliminaries}

\chapter{Membrane computing}

Membranes intro

Natural computing is a recent field of research wchi tries to imitate nature in the way it "computes", learning new computing models and computing paradigms experimented for billions of years by nature.

Neural networks, genetic algorithms and DNA computing are already well established research fields.

However, nature computes not only at the neural or genetic level, but also at the cellular level. In general, any non-trivial biological system has a hierarchical structure where objects and information flows between regions, what can be interpreted as a computation process.

% The notion of membrane

The regions are typically delimited by various types of membranes at different levels from cell membranes, through skin membrane to virtual membranes which delimits different parts of an ecosystem.
This hierarchical system can be seen in other field such as distributed computing, where again well delimited computing units coexist and are hierarchically arranged in complex systems from single processors to the internet.

Membranes keep together certain chemicals or information and selectively determines which of them may pass through.

% The notion of membrane structure

From these observations, Paun \cite{Paun98} introduces the notion of a membrane structure as a mathematical representation of hierarchical architectures composed of membranes. It is usually represented as a Venn diagram with all the considered sets being subsets of a unique set and not allowed to be intersected. Every two sets are either one the subset of the other, or disjoint.






P systems

In previous chapter we introduced the notions of membrane and membrane structure.

% Place objects in the regions.

The next step is to place certain objects in the regions delimited by the membranes. The objects are identified by their names, mathematically symbols from a given alphabet.

% Multisets of objects.

Several copies of the same object can appear in a region, so we will work with multisets of objects.

% Evolution rules

In order to obtain a computing device, we will allow the objects to evolve according to evolution rules. Any object, alone or together with another objects, can be transformed in other objects, can pass through a membrane, and can dissolve the membrane in which it is placed.

% Parallelism

All objects evolve at the same time, in parallel manner across all membranes.

% Priorities

The evolution rules are hierarchizes by a priority relation, which is a partial order.

% P system

These aspects all together forms a P system as introduced in \cite{Paun98}.

\chapter{P systems}
The concept of P systems was first introduced in a 1998 report by Gheorghe Păun in the report \cite{Paun98}.

Some text about P systems. Gheorge Paun established this topic in 1998 with book Membrane Computing. There have been lot of research in this field of study.

\section{Definitions}
How is P system defined. Mostly taken from my article. Multiset rewriting. Accepting vs generating model, active or passive membranes, quality vs quantity aspects.

\section{P system variants}
We are interested in computation power of various variants of P systems. Especially those that are universal (Turing completeness).

\subsection{Parallelism options}
Maximal parallelism, minimal parallelism, n-parallelism, sequential models.

\subsection{Contextivity rules}
Context rules vs cooperational rules, catalytic rules, symmetric cooperational rules, catalytic rules, promoters, inhibitors, context-free rules.

\subsection{Priority rules}
\subsection{Energy of membranes}
\subsection{Calculi of Looping Sequences}

\section{Case studies}
Vultures in Pyrenees, Scavangers of Pyrenees.

\chapter*{Conclusions}
\addcontentsline{toc}{chapter}{Conclusions}
We have studied several variants of sequential P systems in order to obtain universality without using maximal parallelism. A variant with rewriting rules that can use inhibitors was shown to be universal in both generating and accepting case. The generating model is able to simulate maximal parallel P system and the accepting model can simulate a register machine.

In addition, we have defined a new notion of vacuum, which is immediately created in the region that becomes empty. A new P system variant was introduced, which allowed the vacuum to be used on the left side of P system rewriting rules. The accepting case of this variant was shown to be universal by direct implementation of a register machine.

A lot of other variants deserve to be combined with the vacuum, so we suggest to research them more
thoroughly(non-cooperative rules, rules with priorities, decaying objects, deterministic steps, . . . ).
\backmatter

% \begin{thebibliography}{1}
\bibliography{diz}
% \end{thebibliography}

\appendix

\chapter{Štatistiky}
Tu budú štatistiky.

\printindex

\end{document}
